\documentclass[a4paper]{scrartcl}
\usepackage[
    typ=kl,
    klausurtyp=klausur,
    fach=Informatik,
    lerngruppe=XX,
    erwartungshorizontAnzeigen,
    loesungen=folgend,
    module={Aufgabenpool}
]{schule}

\date{xx.xx.xxxx}
\title{Beispielaufgabe}

\setAufgabenpoolOptionen{optionA,optionB}

\begin{document}
    \inputOnce{bsp_text}

    \begin{aufgabe}[subtitle=Beispielaufgabe, points=6]
        Ein wirklich anderer Text für eine Aufgabe. Dazu gibt es auch die \autoref{fig:objektdiagramm}

        \begin{figure}[ht]
            \centering
            \includegraphics[width=0.95\linewidth]{\getBasedir/../aufgabe-1}
            \caption{Ein Objektdiagramm}
            \label{fig:objektdiagramm}
        \end{figure}

        \ifAufgabenpoolOptionTF{optionA}{
            Mache etwas damit.
        }{
            Dieses ist total egal.
        }%
        Sei glücklich.%
        \ifAufgabenpoolOptionT{optionB}{
            Aber mache etwas wirklich gutes.
        }

        \ifAufgabenpoolOptionTF{optionA}{
            \erwartung{macht etwas damit}{5}
        }{
            \erwartung{macht etwas anderes damit}{5}
        }
        \erwartung{lässt anderes sein}{1}
    \end{aufgabe}
    \begin{loesung}
        Eigentlich ist es ganz egal.
    \end{loesung}
\end{document}
